\lecture{9}
\title{Unit-Sample Response and Convolution}

\begin{document}

\subsection{Lecture and Recitation: Unit-Sample Response and Convolution}

Let's just consider LTI systems between DT signals for now.

\textbf{Definition.} Given two signals $a[n]$ and $b[n]$,
their \textbf{convolution} is $(a * b)[n] := \sum_{m = -\infty}^{\infty} a[m] \cdot b[n - m]$.

\begin{theorem}{LTI = Convolution}
    Suppose $\delta[n] \mapsto h[n]$ under an LTI system on DT signals.  Then:
    \[ x[n] \mapsto \sum_{m = -\infty}^{\infty} x[m] \cdot h[n - m] =: (x * h)[n]. \]
\end{theorem}

\begin{proof}
    Immediate from linearity and time-invariance
    after writing $x[n] = \sum_{m = -\infty}^{\infty} x[m] \delta[n - m]$.
\end{proof}

\textbf{Definition.} The \textbf{unit-sample response} of a system on DT signals
is the image $h[n]$ of the Kronecker delta signal $\delta[n]$.

It's not hard to generalize this to CT signals, too.

\textbf{Definition.} Given two signals $f(t)$ and $g(t)$,
their \textbf{convolution} is $(f * g)(t) := \int_{-\infty}^{\infty}
f(\tau) \cdot g(t - \tau) \, \mathrm{d}\tau$.

\textbf{Definition.} The \textbf{impulse response} of a system on CT signals
is the image $h(t)$ of the Dirac delta signal $\delta(t)$.

\begin{theorem}{LTI = Convolution}
    Suppose $\delta(t) \mapsto h(t)$ under an LTI system on CT signals.  Then:
    \[ x(t) \mapsto \int_{-\infty}^{\infty} x(\tau) \cdot h(t - \tau) \, \mathrm{d}\tau
    =: (x * h)(t). \]
\end{theorem}

The point is that every LTI system is a convolution,
defined solely by a single unit-sample / impulse response.

\end{document}